\uuid{gIIO}
\titre{Système linéaire/Méthode du Pivot de Gauss}
\auteur{hanani ; mzari}
\organisation{exo7}
\licence{CC BY-NC-SA 4.0}
\datecreate{2020-01-20}
\niveau{L1}
\module{Algèbre}
\chapitre{Déterminant, système linéaire}
\sousChapitre{Système linéaire, rang}
\difficulte{1}
\theme{Systèmes d'équations linéaires, Système linéaire, Méthode du Pivot de Gauss}
\exo7qcmid{lille-2.systemes.003}

\contenu{
\begin{question}
\qclassification{200.03}
\qtags{motcle=Système linéaire/Méthode du Pivot de Gauss}
On considère le système d'équations, d'inconnue $(x,y,z)\in \Rr^3$ :
$$(\mathtt{S})  \; \left\{\begin{array}{rcc}
x-y+2z&=&1\\
-2x+2y-4z&=&-2\\
3x-3y+6z&=&3.\end{array}\right.$$
Quelles sont les assertions vraies ?
\begin{answers}
\good{$(\mathtt{S}) \Leftrightarrow
x-y+2z=1.$ }
\good{L'ensemble des solutions de $(\mathtt{S})$ est un plan}.
\bad{$(\mathtt{S})$ n'admet pas de solution}.
\bad{$(\mathtt{S})$ admet une unique solution}.
\end{answers}
\begin{explanations}
$ (\mathtt{S}) \Leftrightarrow
x-y+2z=0$.
L'ensemble des solutions de $(\mathtt{S})$ est un plan.
\end{explanations}
\end{question}
}
